% ─────────────────────────────────────────────────────────────────────────────
\section{Visualization Package (\texttt{src/visualization/})}
\label{sec:visualization}
% ─────────────────────────────────────────────────────────────────────────────

The visualization package renders the live simulation using SFML.  It is always compiled
in; when the binary is run with \texttt{--headless} the window is simply not created.

\begin{figure}[H]
  \centering
  \includegraphics[width=0.82\textwidth]{visualization_classes}
  \caption{Visualization package class diagram.}
  \label{fig:viz}
\end{figure}

% ─────────────────────────────────────────────────────────────────────────────
\subsection{\texttt{Renderer} (class)}

Stateless drawing primitives for the world view.  Reuses pre-allocated SFML shapes
(\texttt{circle\_}, \texttt{triangle\_}, \texttt{rect\_}) to avoid per-frame allocation.

\begin{table}[H]
\centering\small
\caption{\texttt{Renderer} public interface}
\begin{tabularx}{\textwidth}{@{}lX@{}}
\toprule
\textbf{Method} & \textbf{Description} \\
\midrule
\texttt{draw\_background(target, width, height)} & Fill background with dark color \\
\texttt{draw\_grid(target, width, height, cell\_size)} & Draw spatial-grid overlay (when \texttt{show\_grid}) \\
\texttt{draw\_boundaries(target, width, height)} & Draw world-edge rectangle \\
\texttt{draw\_food(target, food)} & Small green circles for active food \\
\texttt{draw\_agent(target, agent, selected)} & Triangle (predator) or circle (prey) \\
\texttt{draw\_vision\_range(target, agent)} & Dashed circle at \texttt{vision\_range} \\
\texttt{draw\_sensor\_lines(target, agent, agents, food)} & Lines to nearest predator/prey/food \\
\texttt{species\_color(species\_id)} & Deterministic hue (see below); static \\
\bottomrule
\end{tabularx}
\end{table}

\textbf{Public flags:} \texttt{show\_grid}, \texttt{show\_vision}, \texttt{show\_sensors}
(\texttt{bool}); \texttt{dead\_fade\_alpha} (\texttt{float}).

\subsubsection{Species Color Algorithm}

Hue for species $n$:
\[
  H = (n \times 137.508^\circ) \bmod 360^\circ
\]
The golden angle ($137.508^\circ$) ensures adjacent species IDs produce maximally
dissimilar hues.  Saturation = 1.0, Value = 0.9 (HSV).  Converted to RGB via standard
HSV-to-RGB formula.

% ─────────────────────────────────────────────────────────────────────────────
\subsection{\texttt{OverlayStats} (struct)}

All data needed by \texttt{UIOverlay::draw()}.

\begin{table}[H]
\centering\small
\caption{\texttt{OverlayStats} fields}
\begin{tabularx}{\textwidth}{@{}lX@{}}
\toprule
\textbf{Field} & \textbf{Description} \\
\midrule
\texttt{generation}          & Current generation index \\
\texttt{tick}                & Tick within current generation \\
\texttt{alive\_predators}    & Live predator count \\
\texttt{alive\_prey}         & Live prey count \\
\texttt{best\_fitness}       & Best fitness in current generation \\
\texttt{avg\_fitness}        & Mean fitness in current generation \\
\texttt{num\_species}        & Number of active species \\
\texttt{fps}                 & Measured frames per second \\
\texttt{speed\_multiplier}   & Simulation speed multiplier (1–16) \\
\texttt{paused}              & Simulation is paused \\
\texttt{fast\_forward}       & \texttt{[FF]} HUD label when true \\
\texttt{selected\_agent}     & Index of selected agent ($-1$ = none) \\
\texttt{selected\_energy}    & Energy of selected agent \\
\texttt{selected\_age}       & Age of selected agent \\
\texttt{selected\_kills}     & Kill count of selected agent \\
\texttt{selected\_food\_eaten}& Food eaten by selected agent \\
\texttt{selected\_fitness}   & Fitness of selected agent \\
\texttt{selected\_genome\_complexity} & Node + connection count of selected genome \\
\bottomrule
\end{tabularx}
\end{table}

% ─────────────────────────────────────────────────────────────────────────────
\subsection{\texttt{UIOverlay} (class)}

Renders HUD panels, fitness chart, and neural network visualization.

\begin{table}[H]
\centering\small
\caption{\texttt{UIOverlay} public interface}
\begin{tabularx}{\textwidth}{@{}lX@{}}
\toprule
\textbf{Method} & \textbf{Description} \\
\midrule
\texttt{initialize(font\_path)} & Load font from path (or fallback); return success \\
\texttt{draw(target, stats, genome*)} & Render all overlay elements \\
\texttt{has\_font()} & Whether font loaded successfully \\
\texttt{push\_fitness(best, avg)} & Append to rolling history deques \\
\texttt{set\_activations(vals)} & Store \{node\_id → activation\} for NN panel \\
\bottomrule
\end{tabularx}
\end{table}

\textbf{Private members:}
\texttt{best\_history\_} and \texttt{avg\_history\_}: \texttt{deque<float>} capped at
\texttt{CHART\_MAX\_POINTS = 150}.
\texttt{node\_activations\_}: \texttt{unordered\_map<uint32\_t, float>}.
\texttt{font\_}: \texttt{sf::Font}; \texttt{font\_loaded\_}: \texttt{bool}.

\subsubsection{\texttt{draw\_nn\_panel()} Algorithm}

Layout:
\begin{enumerate}[nosep]
  \item Arrange input nodes on the left column, output nodes on the right, hidden nodes
    in intermediate columns ordered by their position in \texttt{eval\_order}.
  \item For each node, draw a circle coloured by activation value:
    \[
      \begin{aligned}
        v \le -1 &\Rightarrow \text{blue}(30,30,200); \\
        v = 0   &\Rightarrow \text{gray}(180,180,180); \\
        v \ge 1 &\Rightarrow \text{orange}(220,120,20)
      \end{aligned}
    \]
    Intermediate values interpolated via \texttt{std::clamp}.
  \item Draw enabled connections as lines; weight modulates line thickness.
\end{enumerate}

\subsubsection{\texttt{draw\_fitness\_chart()} Algorithm}

Renders two polylines over the rightmost panel area: \texttt{best\_history\_} in orange
and \texttt{avg\_history\_} in yellow.  Y-axis auto-scales to
$[\min(\mathit{history}), \max(\mathit{history})]$.

% ─────────────────────────────────────────────────────────────────────────────
\subsection{\texttt{VisualizationManager} (class)}

Owns the SFML window, camera, input handling, and coordinates rendering.

\begin{table}[H]
\centering\small
\caption{\texttt{VisualizationManager} private members}
\begin{tabularx}{\textwidth}{@{}lX@{}}
\toprule
\textbf{Member} & \textbf{Description} \\
\midrule
\texttt{config\_}           & Run parameters \\
\texttt{window\_}           & SFML window (null in headless) \\
\texttt{camera\_view\_}     & Zoomable/pannable world camera \\
\texttt{renderer\_}         & World drawing primitives \\
\texttt{overlay\_}          & HUD overlay \\
\texttt{overlay\_stats\_}   & Updated each \texttt{render()} call \\
\texttt{paused\_}           & Simulation pause state \\
\texttt{fast\_forward\_}    & FF mode (skip rendering) \\
\texttt{speed\_multiplier\_}& Ticks per rendered frame (1–16) \\
\texttt{selected\_agent\_}  & Index of clicked agent ($-1$) \\
\texttt{zoom\_level\_}      & Camera zoom $\in [0.1, 10.0]$ \\
\texttt{dragging\_}         & Right-click pan active \\
\bottomrule
\end{tabularx}
\end{table}

\begin{table}[H]
\centering\small
\caption{\texttt{VisualizationManager} public interface}
\begin{tabularx}{\textwidth}{@{}lX@{}}
\toprule
\textbf{Method} & \textbf{Description} \\
\midrule
\texttt{initialize()} & Create SFML window; load font \\
\texttt{render(sim, evolution)} & Full frame: clear → draw world → draw HUD → display \\
\texttt{should\_close()} & Window closed or Escape pressed \\
\texttt{handle\_events()} & Process SFML event queue \\
\texttt{is\_paused(), speed\_multiplier()} & Input state queries \\
\texttt{should\_reset() / clear\_reset()} & Reset-requested flag \\
\texttt{should\_step() / clear\_step()} & Single-step flag \\
\texttt{is\_fast\_forward() / clear\_fast\_forward()} & FF flag \\
\texttt{selected\_agent()} & Selected agent index \\
\texttt{set\_selected\_activations(vals, idx\_map)} & Invert \texttt{node\_index\_map} → \{id→activation\}; forward to overlay \\
\texttt{push\_fitness(best, avg)} & Delegate to \texttt{overlay\_.push\_fitness()} \\
\bottomrule
\end{tabularx}
\end{table}

\subsubsection{Key Bindings}

\begin{table}[H]
\centering\small
\caption{Keyboard/mouse event map}
\begin{tabular}{@{}lll@{}}
\toprule
\textbf{Input} & \textbf{Action} & \textbf{Flag set} \\
\midrule
\texttt{Space}        & Toggle pause                        & \texttt{paused\_} \\
\texttt{R}            & Request simulation reset            & \texttt{reset\_requested\_} \\
\texttt{N} (step)     & Advance one tick while paused       & \texttt{step\_requested\_} \\
\texttt{H}            & Toggle fast-forward mode            & \texttt{fast\_forward\_} \\
\texttt{+} / \texttt{-}& Increase/decrease speed multiplier & \texttt{speed\_multiplier\_} \\
\texttt{G}            & Toggle spatial-grid overlay         & \texttt{renderer\_.show\_grid} \\
\texttt{V}            & Toggle vision-range circles         & \texttt{renderer\_.show\_vision} \\
\texttt{S}            & Toggle sensor lines                 & \texttt{renderer\_.show\_sensors} \\
Mouse wheel           & Zoom camera                         & \texttt{zoom\_level\_} \\
Right-click drag      & Pan camera                          & \texttt{dragging\_} \\
Left-click            & Select nearest agent                & \texttt{selected\_agent\_} \\
\texttt{Escape}       & Close window                        & \texttt{running\_} \\
\bottomrule
\end{tabular}
\end{table}

\subsubsection{Camera Model}

\texttt{zoom\_level\_} is clamped to $[0.1, 10.0]$.  The \texttt{sf::View} is scaled
reciprocally: a zoom of 2.0 halves the viewport area.  Right-click drag stores the
cursor position at button-press (\texttt{drag\_start\_}) and updates the view center by
the delta each \texttt{MouseMoved} event.
